\documentclass[11pt]{article}
\usepackage{amsmath,amssymb,amsthm,hyperref}
\newtheorem{lemma}{Lemma}\newtheorem{proposition}{Proposition}
\begin{document}
\title{The different of the torsion cover and a logarithmic refinement}
\author{}\date{9 November 2026}\maketitle

\begin{lemma}
Let $V$ be an excellent discrete valuation ring of characteristic zero,
and let $E/V$ be an elliptic scheme. For a positive integer $n$ write
$E[n]=\operatorname{Spec} A$ and let $\widetilde A$ be its normalization.
Then the trace complementary modules satisfy
\[
 A^\vee=n^{-1}A,\qquad
 \widetilde A^\vee\subseteq n^{-1}\widetilde A.
\]
Consequently the different of every field component of
$\widetilde A/V$ contains $n$. The same containment holds for every
intermediate field of any such component.
\end{lemma}
\begin{proof}
The multiplication morphism $[n]:E\to E$ is finite flat and is a
local complete intersection. On invariant differentials its derivative
is multiplication by $n$. The different equals the K\"ahler different
for finite flat local complete intersections; the smooth-source formula
therefore makes the different of $[n]$ the ideal $(n)$.
These facts are Stacks Project, Tag 0BWB, Lemmas 49.12.3 and 49.12.6.
Base change along the zero section gives $E[n]/V$. Both the module of
relative differentials and its zeroth Fitting ideal commute with this
base change, so its different is also $(n)$. Its relative dualizing
module is invertible. Under the generic trace pairing this is the
complementary module $A^\vee$; the different is its inverse, whence
$A^\vee=n^{-1}A$. This is also immediate locally from the Jacobian
description of the trace element in Tag 0BWB, Lemma 49.12.4, and the
trace identification in Tag 0BW0, Lemmas 49.8.2--3.

Since $A\subseteq\widetilde A$, trace duality reverses the inclusion:
\[
 \widetilde A^\vee\subseteq A^\vee=n^{-1}A
                         \subseteq n^{-1}\widetilde A.
\]
Thus $n\widetilde A^\vee\subseteq\widetilde A$, exactly the stated
different containment. For a tower of normal discrete valuation rings
$V\subseteq W\subseteq U$, the different tower formula gives
$\mathfrak D_{U/V}=\mathfrak D_{U/W}\mathfrak D_{W/V}U$.
The nonnegative valuation of the first factor shows that the upper
bound $v_U(\mathfrak D_{U/V})\le v_U(n)$ descends to
$v_W(\mathfrak D_{W/V})\le v_W(n)$. This proves the last assertion.
All trace arguments apply to each component after localization; no
perfectness of the residue field is needed.
\end{proof}

\begin{proposition}
Let $X_H=X_1(N)/H$ with $-1\in H$, let $d$ be the degree of its
$t=j/1728$ map, and let $g$ be its genus. With the stable Faltings
height and the determinant $L^2$ metric used in the evaluation theorem,
\[
 h_F(\operatorname{Jac}X_H)
 \le\frac d2\bigl(\log12+\log N+2\log d\bigr)
            +g\log(2/\sqrt\pi).
\]
\end{proposition}
\begin{proof}
Use the same auxiliary splitting and semistability field $K$ as in
the evaluation theorem. At a prime $\mathfrak p$ of $K$ over $p>3$,
localize the $t$-line at the generic point of the special fiber.
The family
\[
 E_t:y^2=x^3-\frac{3t}{t-1}x-\frac{2t}{t-1}
\]
is an elliptic scheme over this discrete valuation ring, because its
discriminant is $1728t^2/(t-1)^3$, a unit there. The generic primitive
$N$-torsion cover modulo $H$ is exactly the function field cover of
$X_H$, as proved in the integral-model lemma of the evaluation paper.
It is a quotient of a union of field components of $E_t[N]$.
The preceding lemma bounds the different at every place above
$\mathfrak p$ by its valuation of $N$. Thus the vertical discriminant
coefficient is at most
\[
            d\,v_{\mathfrak p}(N).
\]
Summing and dividing by $[K:\mathbb Q]$ gives at most $d\log N$
from all $p>3$. At $p=2,3$ the generic Lenstra bound used in the
evaluation theorem contributes at most
$d(\log6+2\log d)$. The horizontal branch contribution to the
fiber at $t=2$ is at most $d\log2$. Hence that fiber order has
normalized logarithmic discriminant at most
$d(\log12+\log N+2\log d)$. The proved differential-evaluation
inequality now yields the result.
\end{proof}

\paragraph{Consequence for the retained modular realization.}
Retain $R=B_*M$, $D_0=4608\cdot210^{300}$ and
$B_*=\prod_{p\le76176}p$. The source bounds are
$N\le210^{300}R^2$, $d<D_0R^3$, and $g\le d$.
The exact integer inequalities $D_0^2<B_*$ and $144<B_*$ imply,
with $\ell=\log R$,
\[
 \log N\le\tfrac52\ell,\qquad
 \log d\le\tfrac72\ell,\qquad
 \log12\le\tfrac12\ell.
\]
Thus $h_F(J_H)\le6d\ell$. The identical fixed-dimensional isogeny
calculation, with $8D_0R^3\ell$ in place of
$8D_0R^3\ell^2$, gives $h_F(A)\le9D_0R^3\ell$ and
\[
 h(a)\le19D_0R^3\ell
       \le C M^3(1+\log M),\qquad
 C=19D_0B_*^3\log B_*<2^{330770}.
\]
The final bit bound uses $19<2^5$, $D_0<2^{2413}$,
$B_*<2^{109445}$, and $\log B_*<2^{17}$.

\paragraph{Potential scalar use.}
A fourth-power scalar norm has strict coefficient orders
$v(Z_0)>23\nu/8$ and $v(Z_1)>19\nu/8$. Hence
$M^3\ell_0^{L_N}\mid c$ with $\ell_0=80989$,
$L_N=\lfloor19(2N-1)/8\rfloor+1$, and $\ell_0\nmid M$.
The new estimate then yields
\[
 h(\beta)\le C\ell_0^{-L_N}c(1+\tfrac13\log c).
\]
The sharper Lucas growth estimate of order $c\log c$ would absorb
this term for every sufficiently large fixed $N$, allowing the
$R^8$ coordinate to be deleted. This last application is under
separate audit.
\end{document}
